\section{Three rows}\label{sec:three}

Throughout, $m\ge3$, the rows are $1,2,3$ from the top, and the tilted rectangles are the
diamonds $D(i,1)$ (Lemma~\ref{lem:tilted-narrow}). Theorem~\ref{thm:rowbound} and
Lemma~\ref{lem:cross} give $f(3,m)=m+2$, a result also proved by the Omega Institute~\cite{trureturing2026}. The columns are $a=\{1\}$, $b=\{2\}$, $c=\{3\}$,
$A=\{1,2\}$, $B=\{1,3\}$, $C=\{2,3\}$ and $X=\{1,2,3\}$, with $a$ and $c$ \emph{outer}. A
blackout without empty columns is its \emph{word} $B_1\cdots B_m$, and each pair of $F$
(Section~\ref{sec:conv}) is named by its double letter.

\begin{proposition}[the two types]\label{prop:threeclass}
Let $S$ be a blackout of the $3\times m$ grid with $m+2$ points.
\begin{enumerate}[(a)]
\item If $S$ is valid, its word is one $X$ among singletons (type T) or two different doubles
among singletons (type P).
\item For every word of type T or P, $G_S$ is a tree, so axis-aligned rectangles have
distinct presentations.
\item There are $m3^{m-1}+\binom m2\cdot6\cdot3^{m-2}=m^23^{m-1}$ words of type T or P.
\end{enumerate}
\end{proposition}

\begin{proof}
(a) By Lemmas~\ref{lem:H} and~\ref{lem:mantel}, $|S|=t+e(H)$ with $t\le m$ nonempty columns and
$e(H)\le2$. So every column is nonempty and $\sum_i(|B_i|-1)=2$, and two equal doubles would form
a four-cycle in $G_S$, against Lemma~\ref{lem:girth}. (b) The column $X$, or the two doubles,
connect the three rows, and every other column hangs on one row. (c) Place $X$, or two
different doubles, and choose a singleton for each other column.
\end{proof}

These words are the \emph{candidates}, in the class $X$ and the classes $AC$, $CA$, $AB$,
$BA$, $BC$ and $CB$ named by their doubles from left to right. Let $AC_m$ hold the valid words
of width $m$ in $AC$, and likewise for the others.

\begin{lemma}[diamond test]\label{lem:diamond}
A candidate $S$ with pair set $F$ is valid if and only if every diamond $D(i,1)$ passes the
test below, where $P=D(i,1)\setminus S$.
\begin{enumerate}[(a)]
\item If $|P|\ge3$ or $P$ is empty, the diamond passes.
\item If $P=\{(i+1,1),(i+1,3)\}$, the diamond fails if and only if $B\in F$.
\item If $P=\{(i,2),(i+2,2)\}$, it fails if and only if $B_i\cap B_{i+2}$ meets $\{1,3\}$.
\item If $P$ is an oblique pair $\{(x,r),(z,s)\}$, it fails if and only if $(x,s),(z,r)\in S$.
\item If $P=\{(x,r)\}$, it fails if and only if $\{r,s\}\in F$ for some $s\in B_x$.
\end{enumerate}
\end{lemma}

\begin{proof}
Axis-aligned rectangles are told apart (Proposition~\ref{prop:threeclass}(b)), three points fix
a rectangle (Lemma~\ref{lem:threecorners}), and an empty diamond has no axis-aligned rival, as
$G_S$ has no four-cycle. Tests (b) to (e) are Lemma~\ref{lem:rival}, a rival needing $B\in F$
in (b) and $\{r,s\}\in F$ in (e). Two diamonds share at most a tip, so a common presentation
would need both middle columns black in rows $1$ and $3$, which no candidate has.
\end{proof}

A \emph{window} is three consecutive letters, at columns $i$ to $i+2$, and \emph{passes} or
\emph{fails} with $D(i,1)$. Its \emph{tips} $(i,2)$ and $(i+2,2)$ are black
exactly when their letters are $b$, $A$, $C$ or $X$.

\begin{lemma}[singleton dichotomy]\label{lem:dichotomy}
In a candidate, a window of three singletons fails if and only if $B\in F$ and it is
$bab$, $bbb$ or $bcb$.
\end{lemma}

\begin{proof}
With both tips black, the diamond shows the vertical pair or one point of the middle column, and
tests (b) and (e) fail it exactly when $B\in F$. With one black tip it shows three corners or an
oblique pair passing test (d), as $(i+1,2)$ is white. With no black tip at least three corners
show.
\end{proof}

\begin{lemma}[six of nine in the class $\{A,C\}$]\label{lem:sixofnine}
Let $S$ be a candidate with $F=\{A,C\}$, and let $s,t$ be singletons.
\begin{enumerate}[(a)]
\item A window $sAt$ fails if and only if $st$ is $bb$, $bc$ or $cb$, also for $t=C$ read
as $t=b$.
\item A window $sCt$ fails if and only if $st$ is $bb$, $ba$ or $ab$.
\item A window with one double, at a tip, passes.
\end{enumerate}
\end{lemma}

\begin{proof}
(a) Column $i+1$ shows only $(i+1,3)$, and the tip $(i+2,2)$ is black for $t\in\{b,C\}$.
With both tips black, test (e) fails, as $2\in B_{i+1}$ and $C\in F$. With only $(i+2,2)$
black, the diamond shows $(i,2)$ and $(i+1,3)$, and test (d) fails it exactly when $s=c$. The case of $(i,2)$ alone is symmetric, and with no black tip three
corners show. Reflection in the middle row gives (b). (c) For $Ast$ the proof of
Lemma~\ref{lem:dichotomy} applies with the tip $(i,2)$ in place of $s=b$, as $B\notin F$.
Reflection and reversal give the other windows.
\end{proof}

By Lemma~\ref{lem:diamond}, the set $AC_m$ holds the words $UAVCW$ of width $m$, with $U$, $V$
and $W$ singleton strings, whose windows pass. Put
$G_{AC}(z)=\sum_{m\ge2}|AC_m|z^m$, define $G_{AB}$ likewise from the words $UAVBW$, and let
$G_X(z)=\sum_{m\ge1}|X_m|z^m$ come from the words $UXW$. With $R=1-3z$, nonempty singleton
strings with a fixed letter next to a double have generating function $z/R$.

\begin{proposition}[the class $\{A,C\}$]\label{prop:threeGAC}
The generating function of $AC$ is
\begin{equation}\label{eq:GAC}
G_{AC}(z)=\frac{z^2-4z^3+10z^4-8z^5-3z^6}{(1-3z)^3},
\end{equation}
and
\begin{equation}\label{eq:GACcoef}
|AC_m|=2(m-1)(3m+2)3^{m-5}\qquad(m\ge4).
\end{equation}
\end{proposition}

\begin{proof}
By Lemma~\ref{lem:sixofnine}(a) and the half turn, a window $tAC$ passes only for $t=a$ and
$ACt$ only for $t=c$, and $AtC$ passes by tests (b) and (e). With Lemmas~\ref{lem:dichotomy}
and~\ref{lem:sixofnine}(c), only windows with a double in the middle can fail. If $V$ is
empty, a nonempty $U$ ends in $a$ and a nonempty $W$ begins with $c$, giving $z^2(1+z/R)^2$.
If $V=t$, Lemma~\ref{lem:sixofnine}(a),(b) leaves $p_t=3,1,2$ last letters for $U$ and
$q_t=2,1,3$ first letters for $W$, for $t=a,b,c$. As $\sum p_t=\sum q_t=6$ and $\sum
p_tq_t=13$, this gives $z^3\sum_t(1+p_tz/R)(1+q_tz/R)=z^3(3+12z/R+13z^2/R^2)$. If
$|V|\ge2$, each end gives $3+6z/R$, for $v_1$ alone or one of the six passing pairs
$(u,v_1)$, and the middle $1/R$, for $9z^4(1-z)^2/R^3$. The three cases sum
to~\eqref{eq:GAC}, and $[z^m]z^kR^{-3}=\binom{m-k+2}23^{m-k}$ gives~\eqref{eq:GACcoef}.
\end{proof}

\begin{theorem}[the number of three-row optima]\label{thm:threecount}
Let $E_m$ count the optima of the $3\times m$ grid with a column black in rows $1$ and $3$.
For $m\ge4$,
\begin{equation}\label{eq:threecount}
a_3(m)=4(m-1)(3m+2)3^{m-5}+E_m.
\end{equation}
\end{theorem}

\begin{proof}
Only $B$ and $X$ are black in rows $1$ and $3$, so by
Proposition~\ref{prop:threeclass}(a) these optima are those outside $AC$ and $CA$. Column reversal maps
$AC$, $AB$ and $BC$ onto $CA$, $BA$ and $CB$, and the reflection in the middle row maps $AB$
onto $CB$, so
\begin{equation}\label{eq:classsum}
a_3(m)=2|AC_m|+4|AB_m|+|X_m|=2|AC_m|+E_m\qquad(m\ge3),
\end{equation}
and~\eqref{eq:GACcoef} gives~\eqref{eq:threecount}.
\end{proof}

\subsection{The generating function}

Put $L=1-2z-2z^2$, $N=1+2z^2$ and $D=LN$.

\begin{theorem}[three-row generating function, parts (b) and (c) computer-assisted]\label{thm:threegf}
\leavevmode
\begin{enumerate}[(a)]
\item The series $\sum_{m\ge3}a_3(m)z^m$ equals $2G_{AC}+4G_{AB}+G_X-z-12z^2$, with
$G_{AB}$ and $G_X$ from Propositions~\ref{prop:GAB} and~\ref{prop:GX}, of denominators
dividing $D^3$ and $D^2$.
\item Its reduced denominator is exactly $Q=R^3L^3N^3$.
\item With
\begin{align*}
P(z)={}&z-3z^2-34z^3+228z^4-641z^5+1307z^6-2186z^7\\
&+2056z^8-236z^9-1060z^{10}+2584z^{11}-4264z^{12}\\
&-1440z^{13}+704z^{14}-1024z^{15}+2080z^{16}+2176z^{17}+384z^{18},
\end{align*}
we have
\begin{equation}\label{eq:PQ}
\sum_{m\ge3}a_3(m)z^m=\frac{P(z)}{Q(z)}-z-12z^2.
\end{equation}
\end{enumerate}
\end{theorem}

\begin{proof}
(a) Use~\eqref{eq:classsum}, with $z+12z^2$ removing widths one and two. (b) See Section~\ref{app:threegf}. (c) Exact computer algebra in SymPy~\cite{sympy2017} gives $P$,
and an independent transfer matrix confirms~\eqref{eq:PQ} exactly through width $100$ (\nolinkurl{three_row_closed_form_check.py}).
\end{proof}

Let $\rho=1+\sqrt3$ and $\sigma=1-\sqrt3$.

\begin{theorem}[three-row closed form, computer-assisted exact algebra]\label{thm:threeclosed}
For every $m\ge3$,
\begin{equation}\label{eq:Eclosed}
E_m=p_+(m)\rho^m+p_-(m)\sigma^m+T_m,
\end{equation}
where $432\,p_\pm(m)=(19\pm11\sqrt3)m^2+(61\pm41\sqrt3)m+104\pm28\sqrt3$ and $T_m=(-2)^{\lfloor m/2\rfloor}\,\kappa(m)$ with $\kappa(m)=(2m^2+41m-104)/216$ for even $m$ and
$\kappa(m)=(7m^2-32m-175)/432$ for odd $m$. Hence $E_m\sim\frac{19+11\sqrt3}{432}m^2\rho^m$. Without the computer, we have $E_m>0$ and
\begin{equation}\label{eq:asymp1}
a_3(m)=\frac4{81}m^23^m-\frac4{243}(m+2)3^m+O(m^2\rho^m).
\end{equation}
\end{theorem}

\begin{proof}
By~\eqref{eq:classsum}, the number $E_m$ is the coefficient of $z^m$ in $4G_{AB}+G_X$, whose
denominator divides $L^3N^3$. Its poles $1/\rho$, $1/\sigma$ and $\pm1/(i\sqrt2)$ are at most
triple, so beyond a polynomial part of degree one, the sequence has the form~\eqref{eq:Eclosed}.
Exact computer algebra gives the twelve rational coefficients. Both sides satisfy the
order-$12$ recurrence of $L^3N^3$ from $m=14$ on and agree for $m=2,\dots,13$, hence always. As $|\sigma|<1$
and $|T_m|=O(m^22^{m/2})$, the asymptotic of $E_m$ follows. Without the computer, the pole
$1/\rho$ of least modulus gives $E_m=O(m^2\rho^m)$, hence~\eqref{eq:asymp1}
by~\eqref{eq:threecount}, and the cross of Lemma~\ref{lem:cross}, with word $Xa\cdots a$,
gives $E_m>0$.
\end{proof}
